\section{狭义相对论的时空观}
\label{sec:狭义相对论的时空观}
狭义相对论的基础是相对性原理和光速不变原理，在此基础上我们导出了洛伦兹变换，那么这种变换会导致什么？本节，我们就将通过洛伦兹变换分析几种特定情形，从而认识狭义相对论观点下的时空到底具有何种形态。\empx{注意！抛弃你的直觉！这已经不再是你熟悉的世界了！}

\subsection{同时的相对性}
\label{subsec:同时的相对性}
狭义相对论最直接的结果就是\uwave{同时的相对性}，这就是说，同时发生的两项事件在另外一个参照系未必是同时发生的，或者说，同时不再具有绝对意义。这里有一个例子，如\xref{fig:同时的相对性}所示，设想两个矩形$S,S'$是两列车厢，当然也是两个参照系，其中车厢$S$静止而车厢$S'$以速度$u$运行在铁轨上。在该背景下，设想在某一时刻，在运动车厢$S'$的中心点$P'$处发出一个闪光
\begin{itemize}
    \item 以运动车厢$S'$为参照系，显然光应同时到达$A'$和$B'$。
    \item 以静止车厢$S$为参照系，根据光速不变，由$P'$发出的光仍然以光速$c$向两侧传播，但是由于车厢在运动，$A'$以$u$接近光，$B'$以$u$远离光，因此$A'$将先于$B'$看到光。
\end{itemize}

\begin{OldFigure}[同时的相对性]
    \inputfigure[scale=0.9]{Chapter14C_Figure01}
\end{OldFigure}
这里我们就看到，在$S'$系中认为同时的事件“$A'$和$B'$同时看到闪光”，在$S$系中却不再同时了“$A'$先于$B'$看到闪光”，这就是同时的相对性。这个结果也很容易可以从洛伦兹变换中得到，设惯性系$S'$中发生了两个事件$(x_1',t_1')$和$(x_2',t_2')$，根据\fancyref{fml:洛伦兹变换}
\gdef\prefixeq{同时的相对性}
\begin{Equation}&[1]
    t_2=\gamma\qty(t_2'+\frac{ux_2'}{c^2})\qquad
    t_1=\gamma\qty(t_1'+\frac{ux_1'}{c^2})
\end{Equation}
记$\delt{t'}=t_2'-t_1'$和$\delt{x'}=x_2'-x_1'$，那么在$S$系中观察到的时间间隔$\delt{t}$是
\begin{Equation}&[2]
    \delt{t}=t_2-t_1=\gamma\qty(\delt{t'}+\frac{u}{c^2}\delt{x'})
\end{Equation}
这就可以解释前面列车上的现象了，这里的两个事件是“$A$看到闪光”和“$B$看到闪光”，在运动车厢$S'$上是同时发生（$\delt{t'}=0$），但并不在同一处发生即（$\delt{x'}\neq 0$），根据\xrefpeq{2}给出的结果，这就很自然的说明这两个事件在静止车厢$S$上不是同时发生，因为$\delt{t}\neq 0$的。这里我们也可以得到一个结论，\empx{只有同时同地发生的两个事件，在其他参照系中才是同时发生的}。

这里也可以写出\xrefpeq{2}相反的变换
\begin{Equation}&[3]
    \delt{t'}=\gamma\qty(\delt{t}-\frac{u}{c^2}\delt{x})
\end{Equation}
同时的相对性可能会给我们带来这样一种迷茫，即，\empx{时间的混乱}，事情发生的前后顺序难道就没有意义了？事情其实没有那么严重，尽管同时具有相对性，但因果律下事物的时序仍然具有绝对意义。例如在$S$系中，原因$(x_1,t_1)$在前，结果$(x_2,t_2)$在后，即$\delt{t}=t_2-t_1>0$，但是狭义相对论的结论之一就是速度$v\leq c$，这就意味着具有因果联系的两个事件间的时间间隔必须满足$\delt{t}\geq \delt{x}/c$，因为结果$(x_2,t_2)$要发生，那么至少要等到原因$(x_1,t_1)$以光速跨越两者间$\delt{x}$的距离才能发生，否则就违背了任何物质的运动速度不能超越光速的结论，这样，既然
\begin{Equation}&[4]
    \delt{x}\leq c\delt{t}
\end{Equation}
注意到\xrefpeq{3}可以改写为（最后一步运用$\delt{t}>0$和$1-(u/c)>0$）
\begin{Equation}
    \delt{t'}=\gamma\qty(\delt{t}-\frac{u}{c^2}\delt{x})\geq\gamma\qty(\delt{t}-\frac{u}{c^2}\delt{t})=\gamma\qty(1-\frac{u}{c})\delt{t}>0
\end{Equation}
这就说明在另外一个参照系下必有$\delt{t'}>0$，这就说明任何参照系下因果时序都不变。

\subsection{尺缩效应}
\label{subsec:尺缩效应}
在伽利略变换下，物体的长度并不会随参照系而变，但是，在洛伦兹变换系下，物体的长度就未必是一个无关参照系的常数了，当然，既然长度可变，我们就先需要确定一个“原本”的长度，很合理的，我们定义相对物体静止时测定的长度是其\uwave{固有长度}。那么如何测量运动时的长度呢？很重要的一点是，当我们获得物体两端的时空坐标$(x_1,t_1)$和$(x_2,t_2)$，它们的时间应当是相等$t_1=t_2$的，此时的$\delt{x}=x_2-x_1$才表示某种长度，我们将其称为物体的\uwave{运动长度}。 
\begin{OldBoxFormula}[尺缩效应]
    所谓\uwave{尺缩效应}，是指运动长度$\delt{x}$相较于运动物体的固有长度$\delt{x}'$变短
    \begin{Equation}&[]
        \delt{x}=\gamma^{-1}\delt{x'}
    \end{Equation}
\end{OldBoxFormula}
\begin{Proof}
    证明是容易的，根据\fancyref{fml:洛伦兹变换}中的空间部分（正变换$S\to S'$）
    \begin{Equation}&[1]
        x_1'=\gamma(x_1-ut_1)\qquad x_2'=\gamma(x_2-ut_2)
    \end{Equation}
    将\xrefpeq{1}的中的两式相减
    \begin{Equation}&[2]
        x_2'-x_1'=\gamma\qty\big[(x_2-x_1)-u(t_2-t_1)]
    \end{Equation}
    由于运动长度的测量要求（在$S$系）时间坐标一致，故$t_1=t_2$
    \begin{Equation}&[3]
        x_2'-x_1'=\gamma\qty\big[x_2-x_1]
    \end{Equation}
    记运动长度$\delt{x}=x_2-x_1$和固有长度$\delt{x'}=x_2'-x_1'$，并将$\gamma$移项
    \begin{Equation}*
        \delt{x}'=\gamma\delt{x}\qquad \delt{x}=\gamma^{-1}\delt{x}'\qedhere
    \end{Equation}
\end{Proof}

尺缩效应，或者说，长度的收缩效应纯属时空的性质，长度在观测中确实变短了，但是物质结构层面上没有发生任何的变化，这与热胀冷缩中那种实在的膨胀和收缩是完全不同的事情。

\subsection{钟慢效应}
\label{subsec:钟慢效应}
钟慢效应和尺缩效应是狭义相对论最为出名的两个结论，它们的实质与名称一样是对偶的
\begin{itemize}
    \item 尺缩效应，即运动物体的长度变短，运动物体本身被指定为$S'$系，而我们观察到物体以速度$u$运动的参照系则为$S$系。固有空间距离$\delt{x}'=x_2'-x_1'$是在$S'$系中相对运动物体静止时的空间距离，而我们观测到的运动物体的空间距离$\delt{x}=x_2-x_1$则是$S$系中的观测结果，而尺缩即$\delt{x}<\delt{x}'$。\empx{测量空间距离时要保证时间上它们属于同一时刻}。
    \item 钟慢效应，即运动物体的时间减缓，运动物体本身被指定为$S'$系，而我们观察到物体以速度$u$运动的参照系则为$S$系。固有时间间隔$\delt{t}'\hspace{0.15em}=t_2'\hspace{0.15em}-t_1'\hspace{0.15em}$是在$S'$系中相对运动物体静止时的时间间隔，而我们观测到的运动物体的时间间隔$\delt{t}\hspace{0.15em}=t_2\hspace{0.15em}-t_1\hspace{0.15em}$则是$S$系中的观测结果，而钟慢即$\delt{t}\hspace{0.15em}>\delt{t}'\hspace{0.15em}$。\empx{测量时间间隔时要保证空间上它们属于同一位置}。
\end{itemize}

这里应说明的是，尺缩要求的同一时刻是指测量者所在的$S$系下的同时，即$t_1=t_2$，钟慢所要求的同一位置是指运动物体所在的$S'$系下的同一位置，即$x_1'=x_2'$，两者刚好是相反的，其实也正是这种相反的要求，才导致了“尺缩（减小）”和“钟慢（增大）”两种相反的效应结果。


这里时间的相对性比空间的相对性更难理解，我们说，运动物体的时间减缓，只是当我们处于观测到该物体运动的参照系时得出的结论，运动物体本身的时间流逝并没有变化，这也就是所谓\uwave{固有时间}的想法，\empx{每个参照系都有自己的固有时间，每一个}！至于说，我们观测到的运动物体上的所谓\uwave{运动时间}与运动物体自身的固有时间，则完全是两样不同的东西，并不非得相等。
\begin{OldBoxFormula}[钟慢效应]
    所谓\uwave{钟慢效应}，是指运动时间$\delt{t}$相较于运动物体的固有时间$\delt{t}'$变长
    \begin{Equation}&[]
        \delt{t}=\gamma\delt{t'}
    \end{Equation}
\end{OldBoxFormula}
\begin{Proof}
    证明是容易的，根据\fancyref{fml:洛伦兹变换}中的时间部分（逆变换$S'\to S$）
    \begin{Equation}&[1]
        t_1=\gamma\qty(t_1'+\frac{ux_1'}{c^2})\qquad 
        t_2=\gamma\qty(t_2'+\frac{ux_2'}{c^2})
    \end{Equation}
    将\xrefpeq{1}的中的两式相减
    \begin{Equation}&[2]
        t_2-t_1=\gamma\qty[(t_2'-t_1')-\frac{u}{c^2}(x_2'-x_1')]
    \end{Equation}
    由于运动长度的测量要求（在$S'$系）空间坐标一致，故$x_1'=x_2'$
    \begin{Equation}&[3]
        t_2-t_1=\gamma\qty\big[t_2'-t_1']
    \end{Equation}
    记运动时间$\delt{t}=t_2-t_1$和固有时间$\delt{t'}=t_2'-t_1'$
    \begin{Equation}*
        \delt{t}=\gamma\delt{t}'
    \end{Equation}
    这里结论方向是顺着的，因而不需要像尺缩效应中那样将$\gamma$移项的步骤。
\end{Proof}

钟慢效应导致的一个有趣的结果是，由于我们总会认为运动物体上的时间比我们更慢，不妨设想有一对孪生兄弟，哥哥搭乘亚光速飞船进行太空旅行，弟弟则留在地球上。由于在两者各自的参照系，都认为对方都在高速运动（哥哥认为地球在远离，弟弟认为飞船在远离），这样
\begin{itemize}
    \item 从弟弟的视角，运动的哥哥的时间比自己慢，如果再相见，哥哥应该比自己年轻。
    \item 从哥哥的视角，运动的弟弟的时间比自己慢，如果再相见，弟弟应该比自己年轻。
\end{itemize}
这就出现了矛盾的结果，这就是著名的\uwave{孪生子佯谬}，实际上，它们在各自参照系得到的结论都是正确的，之所以我们觉得矛盾被激化，是因为我们觉得兄弟相见时关于年轻与否的判断必须是唯一确定的，这并没有问题，但是，在狭义相对论适用的范畴内，兄弟实际上是再没有可能相见的，兄弟两人在各自参照系作出的相反判断也就没有关系了。这是因为，如果兄弟两人能再相见，那么哥哥需要在以亚光速离开地球一段时间后再调头返回，而这时，哥哥的飞船在往返过程中就会有加速度，哥哥的飞船相对于地球上的弟弟就不再是惯性系了，涉及加速度的问题的求解需要使用狭义相对论之上的广义相对论，而其结果是，兄弟相见时经历太空旅行的哥哥确实更年轻，因此“佯谬”只是局限于狭义相对论下的结果，正确的称呼是\uwave{孪生子效应}。